\chapter*{Abstract}
\addcontentsline{toc}{chapter}{\protect\numberline{}Abstract}
Single-cell RNA-seq and single-cell ATAC-seq enable cell type atlases across organisms, tissues, and disease. To scale these efforts, lower cost, and support custom multiome assays, we built HyDrop and PUMATAC and systematically quantified method biases on a shared reference.

HyDrop is a flexible, open-source droplet platform for scATAC-seq and scRNA-seq based on dissolvable polyacrylamide hydrogel beads. Beads carry cycled split-pool barcodes and disulfide-cleavable BAC crosslinks. We redesigned the barcode layout for paired-end capture, removed custom priming, dissolved beads in-droplet to overcome diffusion and emulsion instability, tuned primer load to prevent cross-barcoding, and added Pitstop-2 and guanidinium to increase yield. HyDrop-ATAC generated \num{7996} high-quality single-cell profiles from flash-frozen mouse cortex. HyDrop-RNA produced \num{9508} transcriptomes from the same tissue and matched reference scRNA-seq. We further captured a small population of FACS-sorted \textit{Drosophila} neurons, demonstrating applicability to low-input samples. Across mouse cortex, HyDrop-RNA reached a median \num{3389} UMIs and \num{1658} genes per cell, and HyDrop-ATAC recovered a median \num{4148} unique fragments per cell with high TSS enrichment. Chemistry, microfluidic chips, and protocols are open-source and low-cost.

PUMATAC is a universal Nextflow preprocessing framework that unifies trimming, whitelist-aware barcode correction, alignment, and fragmentization across heterogeneous scATAC-seq formats. The pipeline integrates pycisTopic/LSI counting, Otsu-based cell calling, Scrublet doublet removal, Freemuxlet demultiplexing, MACS2 peak calling, and cisTarget motif analysis with containerization and restartable execution, and ships with best-practices code.

We benchmarked eight scATAC-seq methods on human PBMCs: 10x v1, v1.1, v2, 10x Multiome, Bio-Rad ddSEQ, mtscATAC-seq, s3-ATAC, and HyDrop. Across \num{47} libraries and \num{169227} high-quality cells, sequencing efficiency, defined as reads retained in the final cell-by-peak matrix, ranged from \textasciitilde\num{4}% for s3-ATAC and early HyDrop to \textasciitilde\num{28}% for 10x v2. Per-cell sensitivity scaled with efficiency at fixed depth and impacted cell-type annotation, genotype demultiplexing, peak calling, differential accessibility, and transcription factor motif enrichment. Methods also differed in specificity. Promoter and TSS proximity as well as FRIP varied markedly, with ddSEQ more TSS-proximal than 10x, yielding technique-dependent DARs and imperfect cross-technology integration even after batch correction. These results support a two-phase design in practice. Generate consensus peak sets with high-specificity kits, then scale cell numbers with cost-efficient methods at shallower depth.

Together, this work contributes HyDrop, a transparent high-performance droplet assay for scATAC-seq and scRNA-seq, PUMATAC, a reproducible method-agnostic preprocessing stack, and a public PBMC resource of \num{169000}+ scATAC-seq profiles with quantitative guidance on sensitivity, specificity, and depth-versus-cells trade-offs. All designs, code, and datasets are released for reuse.

\chapter*{Layman's abstract}
\addcontentsline{toc}{chapter}{\protect\numberline{}Layman's abstract}
We are living through an era where artificial intelligence is transforming how we interact with the world. While everyday use of AI is still new for most people, biology has already been shaped by it for years. For example, we now have models that predict the 3D structure of proteins, or forecast which genes are switched on in a specific cell at a specific time during development. These models tackle one of the central questions in biology: how the same DNA can give rise to very different cell types, like neurons or muscle cells. Understanding this process is essential for explaining how organisms develop, how cells specialize, and how diseases arise when regulation fails. Accurate machine-learning models let us test such ideas in simulation before turning to the lab, and strongly increase the odds of positive results. However, all models depend on one crucial ingredient: large amounts of high-quality data.

One important source of such data comes from the field of \textit{omics}, which broadly refers to systematic ways of measuring biology at scale. This thesis focuses on two such approaches at the single-cell level: single-cell RNA sequencing (scRNA-seq), which measures which genes are active in each cell, and single-cell ATAC-seq, which measures which region of DNA a cell is "using" at a given time. Together, these methods provide a small part of the puzzle of how cells work.

In the first part of this thesis, we developed HyDrop, an open-source method for generating single-cell data. The method uses tiny droplets of liquid to capture and read out information from thousands of individual cells in parallel. Unlike expensive commercial systems, HyDrop is transparent, flexible, and affordable, making cutting-edge single-cell experiments accessible to many more labs.

In the second part, we systematically compared 8 different methods for single-cell chromatin accessibility sequencing. Different technologies often give very different results, and without careful evaluation it is difficult to know which findings can be trusted. By benchmarking protocols side by side, we showed how trade-offs in cost, sensitivity, and data quality influence the biological insights one can draw. This helps researchers choose the right tool for their question, ensures resources are used efficiently, and improves the reliability of downstream analyses.

Taken together, these contributions meet two essential needs: creating new ways to generate high-quality biological data, and setting clear standards to evaluate the methods that produce it. Our results show that technology choices directly shape scientific conclusions from how cell types are identified to how subtle genetic signals are detected.

Finally, we connect these advances to a broader trend. Just as AI in language and vision achieved breakthroughs by scaling data and computation, biology is entering a similar phase. Sequencing costs are falling, automation is rising, and massive datasets are becoming feasible. With open methods for data generation and rigorous frameworks for evaluation, we can move toward predictive models that link DNA to cellular behavior and ultimately to disease. This thesis contributes to building that foundation.

\chapter*{Vulgariserend abstract}
\addcontentsline{toc}{chapter}{\protect\numberline{}Vulgariserend abstract}

We leven in een tijd waarin kunstmatige intelligentie verandert hoe we met de wereld omgaan. Voor de meeste mensen is dagelijks gebruik van AI nog nieuw, maar wetenschap, waaronder de biologische wetenschappen, wordt er al jaren door gevormd. Er bestaan vandaag de dag modellen die de 3D-structuur van eiwitten accuraat kunnen voorspellen. Tevens bestaan er modellen die voorspellen welke genen in een specifieke cel op een bepaald moment tijdens de ontwikkeling van een embryo actief zijn. Deze modellen focussen op een van de oudste kernvragen in de biologie: hoe kan het dat het DNA van al onze cellen identiek is, maar cellen, zoals neuronen en spiercellen, toch van elkaar verschillen? Het antwoord ligt in genregulatie: verschillende cellen gebruiken verschillende genen. Dit proces verklaart hoe organismen zich ontwikkelen, hoe cellen specialiseren en hoe ziekten ontstaan wanneer regulatie faalt. Nauwkeurige AI-modellen maken het mogelijk ideeën eerst in simulaties te testen voor we ze uittesten, wat de kans op succes vergroot. Alle modellen steunen echter op één cruciale bouwsteen: grote hoeveelheden data van hoge kwaliteit.

Een belangrijke bron daarvan komt uit de zogenaamde "omics"\llap{,} de verzamelnaam voor systematische metingen op grote schaal. Dit proefschrift richt zich op twee benaderingen op single-cell-niveau: single-cell RNA-sequencing (scRNA-seq), dat vastlegt welke genen per cel actief zijn, en single-cell ATAC-seq, dat meet welke regio's van het DNA van een cel op een bepaald moment toegankelijk zijn. Samen leveren deze methoden een stukje van de puzzel van het functioneren van cellen.

In het eerste deel van dit proefschrift hebben we HyDrop ontwikkeld, een open-source methode om single-cell data te genereren. De methode gebruikt kleine druppels vloeistof om informatie uit duizenden individuele cellen parallel te vangen en uit te lezen. In tegenstelling tot dure commerciële systemen is HyDrop open, flexibel en betaalbaar, waardoor geavanceerde single-cell-experimenten voor veel meer laboratoria binnen bereik komen.

In het tweede deel hebben we acht protocollen voor scATAC-seq systematisch vergeleken. Verschillende technologieën leveren vaak uiteenlopende resultaten op en zonder zorgvuldige evaluatie is het lastig te bepalen welke bevindingen betrouwbaar zijn. Door protocollen naast elkaar te benchmarken laten we zien hoe afwegingen in kosten, gevoeligheid en datakwaliteit de biologische inzichten sturen. Dit helpt onderzoekers het juiste instrument voor hun vraag te kiezen, middelen efficiënt in te zetten en de betrouwbaarheid van downstreamanalyses te vergroten.

Samen voorzien deze bijdragen in twee kernbehoeften: nieuwe manieren om hoogwaardige biologische data te genereren en duidelijke standaarden om de methoden die deze data opleveren te beoordelen. Onze resultaten laten zien dat technologiekeuzes wetenschappelijke conclusies direct vormgeven, van de identificatie van celtypen tot het detecteren van subtiele genetische signalen.

Tot slot plaatsen we deze vooruitgang in een bredere ontwikkeling. Net als AI in taal en beeld doorbraken bereikte door data en rekenkracht op te schalen, staat de biologie aan het begin van een vergelijkbare fase. Sequencingkosten dalen, automatisering neemt toe en zeer grote datasets worden haalbaar. Met open methoden voor datageneratie en strikte kaders voor evaluatie kunnen we toewerken naar voorspellende modellen die DNA verbinden met cellulair gedrag en uiteindelijk met ziekte. Dit proefschrift helpt die basis te leggen.